\section{Experimental Evaluation}
\label{sec:evaluation}
% \todo[inline]{R2 on scalability: Note that the scalability is not directly affected by the size of the stream database, $|D|$,
% but rather by the structure of the streams. A small database with complex,
% repetitive patterns can cause combinatorial explosions in candidate queries, leading to runtimes of minutes or hours.
% Conversely, a large database with simple, non-repetitive patterns may be processed in seconds.
% Thus, scalability evaluations based on $|D|$ or the alphabet size $|\Gamma|$ are unreliable, as runtime depends heavily on stream structure.
% One explanation for the difference in the datasets is the size of the queryset.
% If for example, the queryset stays the same after updating the database, then the updated algorithm will terminate quickly
% whereas the baseline algorithm has to rediscover all matching queries. We will add the explanation to the final version of the paper.}

This section presents an evaluation of the algorithms proposed in
\autoref{sec:algos}. We
first introduce our experimental setup (\autoref{sec:exp_setup}), before
comparing our algorithms for incremental event query discovery against a
baseline solution (\autoref{sec:exp_sota}).
% Finally, we evaluate the sensitivity of
% our algorithms with respect to the characteristics of a stream database
% (\autoref{sec:exp_sensitity}), before studying the effectiveness of our
% strategies to guide the application of discovery algorithms
% (\autoref{sec:exp_guidance}).

\subsection{Experimental Setup}
\label{sec:exp_setup}

\sstitle{Datasets}
We use two real-world datasets to evaluate our algorithms, namely
 NASDAQ stock trade events~\cite{eodata}
and the Google cluster traces~\cite{reiss2011google}.
We adopted the methodology for obtaining stream databases from the~\cite{ilminer}.
For each dataset, we define three queries to generate streams, as described below.

The attributes considered in the finance data are:
% The event schema for the finance data comprises three attributes:
stock name, volume, and closing value. The queries to construct streams
capture
two events with stock name `GOOG' (F1); three events with stock names
appearing in the sequence `AAPL', `GOOG', and `MSFT' (F2); and four events
for which the volume remains constant (F3).

Events in the Google cluster traces contain four attributes:
job ID, machine ID, task status, and priority. One query
defines three events running on the same machine with identical job IDs,
with the first and last ones of status `submit', while the
second one has
status `kill' (G1); one query encompasses four events
executing on the same
machine with status `finish' (G2); and one query detects job eviction
resulting from the scheduling of another job on the same
machine (G3).

Using these queries, we generated the stream databases summarized in
\autoref{tab:databases} by evaluating the queries over the
datasets. For each match (of the first at most 1000 matches), we constructed a
stream by including the events preceding the match.
The number of events taken into account for each stream
 is detailed in the column 'stream length'.


\begin{table}[h]
	\caption{Stream databases used in the experiments.}
	\label{tab:databases}
	\small
	% \vspace{-1.2em}
	\centering
	\begin{tabular}{c c c c c }
		\toprule
		\multicolumn{2}{c}{}          & stream
			length& number
			of
			streams\\%& \begin{tabular}[c]{@{}c@{}}support\\
			%threshold\end{tabular} \\
		\midrule
		\multirow{3}{*}{Finance} & F1 & 50                    &
		69-79                                                           \\
		& F2 & 40
		& 990-1000                                                       \\
		& F3 & 25
		& 240-250                                                           \\
		%\midrule
		\multirow{1}{*}{Google}  & G1-G3 & 7              &
		990-1000                                                         \\
		\bottomrule
	\end{tabular}
	%\vspace{-1em}
\end{table}

Note that the runtime is not determined by the size of the stream database,
$|D|$, but rather by the structure of the streams. A small database with intricate and repetitive patterns can trigger
a combinatorial explosion of candidate queries, leading to runtimes of minutes or hours.
Conversely, a large database with simple, non-repetitive patterns may be processed in a matter of seconds.

\sstitle{Baseline}
% \todo[inline]{RS: Add experiments for combined updates.}
We compare the runtime of our algorithms making use of incremental updates against
the baseline approach, in which no prior information is given to the
discovery.


\sstitle{Measures} We primarily measure the efficiency of event query
discovery in terms of the algorithms' runtime. We report averages over five
experimental runs, including error bars (mostly negligible).
%For each stream database, the query discovery was repeated 5 times and the
%average runtime is plotted including the confidence interval error bars.

\sstitle{Environment and Implementation}
We implemented the whole experimental evaluation in Python.
The code and the evaluation routine is publicly available on GitHub\footnote{https://github.com/rebesatt/IncDISCES}.
All experiments were conducted on a server with a E7-4880 CPU @2.5GHz and
1TB of RAM.
\subsection{Baseline Comparison}
\label{sec:exp_sota}

% We first conducted experiments on the datasets described in \autoref{tab:databases}
% by decreasing the space of candidate queries, i.e., by adding streams, by
% increasing the support threshold or both.
We first conducted experiments on the datasets described in \autoref{tab:databases} by reducing
the candidate query space—either by adding streams, increasing the support threshold, or both.

In \autoref{fig:gen_sample}, we plotted the relative time change of the
incremental algorithm
compared to the baseline, when adding up to 10 streams to the different
stream databases.
It can be seen that our incremental algorithm is at least an order of
magnitude faster for all settings.
For the google dataset G1, the improvement is even up to three orders of
magnitude.

In \autoref{fig:gen_supp}, we consider the change of increasing the support
threshold.
The original support threshold is 0.75, which is increased stepwise, up to
1.0.
Here, we see a drop in the relative time change for increasing support
values. This is
especially true for the finance stream databases. One reason for this
observation is the
decreasing search space. For example, the number of queries for F1 and
support threshold
1.0 is four. As consequence, both, the baseline and the incremental
algorithm terminate quickly, in less than a second.

% In \autoref{fig:gen_sampsupp}, we consider the simulatenous change of increasing the support
% threshold and added streams.

In \autoref{fig:gen_sampsupp}, we examine the simultaneous effect
of increasing the support threshold and adding streams.
Like in the previous experiment, the original support threshold is 0.75 and
is increased stepwise up to 1.0. At the same time, for each step, we add
one stream to the stream database.
Similar to \autoref{fig:gen_supp}, we see a drop in the relative time change
for increasing support values, but still the runtime of the incremental algorithm is
always lower than the baseline solution.


% \begin{figure}[t]
% 	\centering
% 	\begin{subfigure}{.63\textwidth}
% 	  \centering
% 	  \includegraphics[width=.82\linewidth]{img/gen_sample.pdf}
% 	  \includegraphics[clip,trim=73em 0em 5.5em 0em,
% 	  width=.14\linewidth]{img/gen_supp.pdf}
% 	  \caption{Adding streams to the stream database}
% 	  \label{fig:gen_sample}
% 	\end{subfigure}\\[1em]%
% 	\begin{subfigure}{.63\textwidth}
% 	  \centering
% 	  \includegraphics[width=\linewidth]{img/gen_supp.pdf}
% 	  \caption{Increasing the support threshold}
% 	  \label{fig:gen_supp}
% 	\end{subfigure}
% 	\begin{subfigure}{.63\textwidth}
% 	  \centering
% 	  \includegraphics[width=\linewidth]{img/gen_sampsupp.pdf}
% 	  \caption{Increasing sample size and the support threshold}
% 	  \label{fig:gen_sampsupp}
% 	\end{subfigure}
% 	\caption{Relative time change of updated algorithm compared to baseline
% 	for different stream databases while decreasing the search space.}
% 	\label{fig:test}
% 	\end{figure}

\begin{figure}[t]
    \centering
    \begin{minipage}[t]{0.48\textwidth} % Left column
        \centering
        \vspace{0pt} % Force alignment at the top
        \begin{subfigure}{\textwidth}
            \centering
            \includegraphics[width=\linewidth]{img/gen_sample.pdf}
            % \includegraphics[clip,trim=73em 0em 5.5em 0em, width=.14\linewidth]{img/gen_supp.pdf}
            \caption{Adding streams to the stream database}
            \label{fig:gen_sample}
        \end{subfigure}\\[1em]
        \begin{subfigure}{\textwidth}
            \centering
            \includegraphics[width=\linewidth]{img/gen_supp.pdf}
            \caption{Increasing the support threshold}
            \label{fig:gen_supp}
        \end{subfigure}
        \begin{subfigure}{\textwidth}
            \centering
            \includegraphics[width=\linewidth]{img/gen_sampsupp.pdf}
            \caption{Increasing sample size and the support threshold}
            \label{fig:gen_sampsupp}
        \end{subfigure}
        \caption{Relative time change of updated algorithm compared to baseline for different stream databases while decreasing the search space.}
        \label{fig:test}
    \end{minipage}
    \hfill
    \begin{minipage}[t]{0.48\textwidth} % Right column
        \centering
        \vspace{0pt} % Force alignment at the top
        \begin{subfigure}{\textwidth}
            \centering
            \includegraphics[width=\linewidth]{img/spec_sample.pdf}
            % \includegraphics[clip,trim=76em 0em 5.5em 0em, width=.14\linewidth]{img/spec_supp.pdf}
            \caption{Deleting streams to the stream database}
            \label{fig:spec_samp}
        \end{subfigure}\\[1em]
        \begin{subfigure}{\textwidth}
            \centering
            \includegraphics[width=\linewidth]{img/spec_supp.pdf}
            \caption{Decreasing support threshold}
            \label{fig:spec_supp}
        \end{subfigure}
        \begin{subfigure}{\textwidth}
            \centering
            \includegraphics[width=\linewidth]{img/spec_sampsupp.pdf}
            \caption{Decreasing sample size and support threshold}
            \label{fig:spec_sampsupp}
        \end{subfigure}
        \caption{Relative time change of updated algorithm compared to baseline for different stream databases while increasing the search space.}
        \label{fig:test2}
    \end{minipage}
\end{figure}



% Next, we evaluated the behaviour of the algorithms when increasing the space
% of candidate queries that need to be searched, by deleting streams from the
% database and/or by decreasing the support threshold.
Next, we evaluated the behaviour of the algorithms when increasing the space
of candidate queries that need to be searched, by deleting streams from the
database and/or by decreasing the support threshold.

In \autoref{fig:spec_samp}, for each dataset, the
incremental
algorithm outperforms the baseline solution. For the google stream database
G2, the
incremental algorithm is twice as fast independently of the number of
streams deleted
from the original database. The highest benefit is observed for the google
dataset G3,  where the updated algorithm is up to 50 times faster.

In \autoref{fig:spec_supp}, we see that for decreasing support threshold for
all datasets, the two algorithms need roughly the same
runtime to discover the queries.
This is because the original setting has a small set of matching queries. As
a consequence, the benefit of the incremental discovery algorithms is also
small, which yields a similar runtime as the baseline
solution.

Similarly, in the experiment shown in \autoref{fig:spec_sampsupp},
where both the support threshold and the number of streams are decreased,
the results follow the same trend. The small set of matching queries in the original
setting limits the advantage of the incremental discovery algorithms, leading to comparable
runtimes with the baseline approach.


% \begin{figure}[t]
% 	\centering
% 	\begin{subfigure}{.63\textwidth}
% 		\centering
% 		\includegraphics[width=.82\linewidth]{img/spec_sample.pdf}
% 		\includegraphics[clip,trim=76em 0em 5.5em 0em,
% width=.14\linewidth]{img/spec_supp.pdf}
% 		\caption{Deleting streams to the stream database}
% 		\label{fig:spec_samp}
% 	\end{subfigure}\\[1em]%
% 	\begin{subfigure}{.63\textwidth}
% 		\centering
% 		\includegraphics[width=\linewidth]{img/spec_supp.pdf}
% 		\caption{Decreasing support threshold}
% 		\label{fig:spec_supp}
% 	\end{subfigure}
% 	\begin{subfigure}{.63\textwidth}
% 		\centering
% 		\includegraphics[width=\linewidth]{img/spec_sampsupp.pdf}
% 		\caption{Decreasing sample size and support threshold}
% 		\label{fig:spec_sampsupp}
% 	\end{subfigure}
% 	\caption{Relative time change of updated algorithm compared to baseline
% 	for different stream databases while increasing the search space.}
% 	\label{fig:test2}
% 	\end{figure}